\section{Resumen Ejecutivo e Introducción}

El presente documento expone la propuesta metodológica refinada para la ejecución del estudio sobre la reinserción e inclusión laboral B2B de personas con antecedentes penales en Chile, con especial foco en las barreras de género que afectan a mujeres anteriormente privadas de libertad.

Entre los años 2025 y 2030, se estima que aproximadamente 30.000 personas recobrarán la libertad en Chile, de las cuales un 8,2\% corresponde a mujeres, caracterizadas por una vulnerabilidad multidimensional agravada, dado que el 85,4\% de ellas son madres y sustentadoras de hogar. \parencite{CaracterizacionPersonasPrivadas2026} Responder a este desafío exige superar la intuición comercial y dotar a la Fundación Juntos por la Reinserción (JxR) de evidencia científica robusta y causalmente válida respecto a los procesos de toma de decisión en el sector empresarial.

La estrategia propuesta integra un diseño metodológico cualitativo-explicativo basado en un estudio de casos múltiples con unidades incrustadas, articulado en tres fases complementarias y un análisis secundario de información administrativa. Esta aproximación garantiza un equilibrio riguroso entre la máxima validez interna y externa y la factibilidad operativa B2B.

A través de esta investigación, la Directiva del Proyecto dispondrá de una matriz de tipología de empresas según su disposición de adopción B2B, una guía práctica de abordaje y prospección comercial para superar puntos de veto organizacionales, recomendaciones para el rediseño de sus tres modelos operativos (Bolsa Laboral, OTEC y CET) y un conjunto de argumentos de valor probados experimentalmente.